\section{Three separated leaves on an odd cycle}\label{sec:threeleaves}

The same equality obstruction resolves the next marking with three roots.
We prove Theorem~\ref{thm:threeleaves} by locating a selected root whose two
cycle neighbors are unselected.

\begin{proof}[Proof of Theorem~\ref{thm:threeleaves}]
Write the three cyclic distances as $d_1,d_2,d_3$. By assumption they are
odd and at least three, and $d_1+d_2+d_3=2k+1$. A minimum cover of $G$
can be chosen to contain all three roots and none of the leaves: replace
any selected leaf whose root is absent by that root, and discard any leaf
whose root is present. Once the roots are selected, the three uncovered
interior paths have $d_i-1$ vertices each. Consequently
\[
 \tau(G)=3+\sum_{i=1}^{3}\frac{d_i-1}{2}=k+2.
\]
There are $n=m=2k+4$ vertices and edges. The star construction in the proof of
Proposition~\ref{prop:positive} gives the upper bound
$m+\tau(G)-1=3k+5$.

The induced cycle has cover number $k+1$. Corollary~\ref{cor:core} therefore
requires $n-1+(k+1)=3k+4$ AND/OR gates. Suppose a circuit has at most this
many total gates, and normalize it. It is positive. The restriction
argument in Corollary~\ref{cor:core} requires at least $k+1$ branches;
the incidence budget now forces exactly that many, each with two consumers. As in the two-leaf proof,
the eligible cycle vertices form a cover $S$ of the cycle and have distinct
first branching carriers with singleton supports. Thus $|S|=k+1$, all
branches are those carriers, and no leaf is eligible.

A minimum cover of an odd cycle has exactly one edge with both endpoints
in $S$. Indeed its complement has $k$ independent vertices. Between
consecutive complementary vertices there are $k$ nonempty strings of
selected vertices containing $k+1$ vertices altogether. Exactly one string
has length two and every other has length one. Denote the unique edge
inside $S$ by $e$.

If a marked root $v\in S$ is not an endpoint of $e$, both its cycle
neighbors are ineligible. Its two cycle edges and its leaf edge then have
three distinct meeting gates consuming $v$: any pair of the other endpoints
includes an ineligible cycle vertex, whose cut carriers have singleton
support by Lemma~\ref{lem:carrier}. But the two exits of $b(v)$ allow at
most two such gates, as proved in Lemma~\ref{lem:slots}. Hence every marked
root in $S$ must be an endpoint of $e$. The roots are pairwise nonadjacent,
so at most one belongs to $S$.

At least one root must belong to $S$. If all three were outside, the path
between two consecutive roots would have $d_i-1$ even interior vertices
and both endpoints unselected. Its interior cover would cost
$1+(d_i-1)/2$. Summing over the three disjoint interiors gives at least
$k+2$ selected vertices, contradicting $|S|=k+1$.
Thus exactly one root $v$ lies in $S$ and is an endpoint of $e$.
The arc between the other two roots which does not pass through $v$ has
odd length and has unselected endpoints. It must contain an edge with both
endpoints selected: otherwise, since no edge can have both endpoints
unselected, selection would alternate along this odd arc and its endpoint
selections would differ. Its selected edge cannot be $e$, which is incident
to $v$ outside the arc. This contradicts uniqueness of $e$.
Therefore $3k+4$ gates are impossible, proving the claimed optimum.
\end{proof}

For the nine-cycle, the three distances are $3,3,3$ and $k=4$, so the
optimum is seventeen. One cover construction uses centers
$0,2,3,5,6,8$. With fresh leaves $9,10,11$ at $0,3,6$, respectively,
the six terms are
\[
 x_0(x_1\OR x_8\OR x_9),\quad x_2(x_1\OR x_3),\quad
 x_3(x_4\OR x_{10}),\quad x_5(x_4\OR x_6),\quad
 x_6(x_7\OR x_{11}),\quad x_7x_8.
\]
The first term costs three gates, the last costs one, and each other term
costs two. Their disjunction uses five further OR gates and expands to
exactly the twelve edge terms, giving the seventeen-gate witness.
